\subsection{ASYMPTOTIC EXPANSION OF A PRIMITIVE}

Let E be a comparison scale on a neighbourhood of \(+\infty\) formed of real functions \(\neq 0\) of constant sign on a neighbourhood of \(+\infty\) ; let f be a regulated vector function defined on an interval \([a, +\infty[\) , with values in a complete normed space E, admitting an asymptotic expansion

\[
\mathbf {f} = \sum_ {\lambda \leqslant \alpha} \mathbf {a} _ {\lambda} g _ {\lambda} + \mathbf {r} _ {\alpha}
\]

to precision \(g_{\alpha}\) relative to E. Suppose further that every primitive \(\int_{a}^{x} g(t) dt\) of a function \(g \in E\) admits an asymptotic expansion with respect to E. In these circumstances we shall see that one can obtain an asymptotic expansion of \(\mathbf{F}(x) = \int_{a}^{x} \mathbf{f}(t) dt\) with respect to E. We distinguish two cases:

\(1^{\circ}\int_{a}^{+\infty}g_{\alpha}(t)dt\) is infinite; then one has \(\int_{a}^{x}\mathbf{r}_{\alpha}(t)dt\ll \int_{a}^{x}g_{\alpha}(t)dt\) (V,p.230, prop. 6); by hypothesis one can obtain an asymptotic expansion of \(\sum_{\lambda \leqslant \alpha}\mathbf{a}_{\lambda}\int_{a}^{x}g_{\lambda}(t)dt\)

to a certain precision \(g_{\rho}\) (V, p.~222); if \(cg_{\sigma}\) is the principal part of \(\int_{a}^{x} g_{\alpha}(t) dt\) one will thus have an asymptotic expansion of \(\int_{a}^{x} \mathbf{f}(t) dt\) to precision \(g_{\min(\rho, \sigma)}\) , with all the terms having indefinitely increasing norms.

\(2^{\circ}\int_{a}^{+\infty}g_{\alpha}(t)dt\) converges; let \(\beta\) then be the smallest of the indices \(\lambda \leqslant \alpha\) such that \(\mathbf{a}_{\lambda}\neq 0\) and such that \(\int_{a}^{+\infty}g_{\lambda}(t)dt\) converges; the integral

\[
\mathbf {C} = \int_ {a} ^ {+ \infty} \left(\mathbf {f} (t) - \sum_ {\lambda <   \beta} \mathbf {a} _ {\lambda} g _ {\lambda} (t)\right) d t
\]

is then convergent, and one can write

\[
\mathbf {F} (x) = \sum_ {\lambda <   \beta} \mathbf {a} _ {\lambda} \int_ {a} ^ {x} g _ {\lambda} (t) d t + \mathbf {C} - \sum_ {\beta \leqslant \lambda \leqslant \alpha} \mathbf {a} _ {\lambda} \int_ {x} ^ {+ \infty} g _ {\lambda} (t) d t - \int_ {x} ^ {+ \infty} \mathbf {r} _ {\alpha} (t) d t.
\]

Then \(\int_{x}^{+\infty}\mathbf{r}_{\alpha}(t)dt\ll \int_{x}^{+\infty}g_{\alpha}(t)dt\) ; if \(cg_{\sigma}\) is the principal part of \(\int_{x}^{+\infty}g_{\alpha}(t)dt\) , and if one has an asymptotic expansion of

\[
\sum_ {\lambda <   \beta} \mathbf {a} _ {\lambda} \int_ {a} ^ {x} g _ {\lambda} (t) d t + \mathbf {C} - \sum_ {\beta \leqslant \lambda \leqslant \alpha} \mathbf {a} _ {\lambda} \int_ {x} ^ {+ \infty} g _ {\lambda} (t) d t
\]

to precision \(g_{\rho}\) , one will as a result have an asymptotic expansion of \(\mathbf{F}\) to precision \(g_{\min(\rho, \sigma)}\) .

So it all amounts to finding asymptotic expansions with respect to E of primitives of functions in E. We have seen how, subject to certain hypotheses on E, prop. 8 of V, p.~233 gives the principal part of such a primitive. Further, the proof of prop. 8 gives the expression for the difference of the two sides of formula (1) (resp. (2)) of V, p.~233, in the form of a primitive of the function \(\frac{1}{|\mu+1|}\left(xf'(x)+f(x)\right)-f(x)\) (resp. \(f(x)g'(x)\) with \(g=f/f'\) ); on forming the principal part of this new primitive, as an asymptotic expansion of the right-hand side of (1) (resp. (2)), one obtains the right-hand side of the sought-for expansion (see V, p.~247-255).

Examples. 1) Let \(f(x) = 1 / \log x\) ( \(x > 1\) ); we have seen that \(\int_{a}^{x} dt / \log t \sim x / \log x\) , and that the difference \(\int_{a}^{x} \frac{dt}{\log t} - \frac{x}{\log x}\) is a primitive of \(1 / (\log x)^2\) ; one can apply prop. 8 again to this function, so obtaining \(\int_{a}^{x} dt / (\log t)^2 \sim x / (\log x)^2\) . By recursion one thus obtains the expansion

\[
\begin{array}{l} \int_ {a} ^ {x} \frac {d t}{\log t} = \frac {x}{\log x} + \frac {x}{(\log x) ^ {2}} \\ \qquad + \frac {2 x}{(\log x) ^ {3}} + \dots + (n - 1)! \frac {x}{(\log x) ^ {n}} + o \left(\frac {x}{(\log x) ^ {n}}\right). \end{array}
\]

Note that, irrespective of n, all the terms of this expansion tend to \(+\infty\) with x.

\begin{enumerate}
\def\labelenumi{\arabic{enumi})}
\setcounter{enumi}{1}
\tightlist
\item
  Let \(f(x) = \frac{e^x}{x^2 + 1}\) ; one can write \(f(x) = \frac{e^x}{x^2} - \frac{e^x}{x^4} + o_1\left(\frac{e^x}{x^4}\right)\) . Prop. 8 gives the expansions
\end{enumerate}

\[
\int_ {a} ^ {x} \frac {e ^ {t}}{t ^ {2}} d t = \frac {e ^ {x}}{x ^ {2}} + \frac {2 e ^ {x}}{x ^ {3}} + \frac {6 e ^ {x}}{x ^ {4}} + o _ {2} \left(\frac {e ^ {x}}{x ^ {4}}\right), \quad \int_ {a} ^ {x} \frac {e ^ {t}}{t ^ {4}} d t = \frac {e ^ {x}}{x ^ {4}} + o _ {3} \left(\frac {e ^ {x}}{x ^ {4}}\right)
\]

whence, on adding,

\[
\int_ {a} ^ {x} \frac {e ^ {t}}{t ^ {2} + 1} d t = \frac {e ^ {x}}{x ^ {2}} + 2 \frac {e ^ {x}}{x ^ {3}} + 5 \frac {e ^ {x}}{x ^ {4}} + o _ {4} \left(\frac {e ^ {x}}{x ^ {4}}\right).
\]

